% Literature Review, section 2. Notes for the whole chapter are in 00_chapter_intro.tex.
\section{Quantifying Coronary Geometry}
\label{sec:lit-measure}

Using coronary geometry as a marker relies on a consistent quantification of vessel shape.
However, clinical studies have quantified coronary geometry inconsistently over the years, and the
chosen measure can determine whether a relationship is even detected at all.

Cui et al.\ (2017) \cite{cui2017bifurcation} measured three angles at the left main bifurcation on
the CCTA of patients referred for invasive angiography.
These were the angle between the LAD and LCX and the angles between the left main artery and each of
its two daughters.
Only the angle between the LAD and LCX was wider in patients with significant stenosis, while the two
angles involving the left main artery did not differ.
A study measuring only the angles involving the left main artery would therefore have found no
relationship between the bifurcation angle and stenosis.

For tortuosity, studies using different measures have disagreed even on the direction of the
association.
Groves et al.\ (2009) \cite{groves2009tortuosity} screened the invasive angiograms of consecutive
patients and graded tortuosity visually as two consecutive turns that reversed the direction of a
major artery.
Patients with such severe tortuosity were less likely to have significant CAD than a randomly chosen
group without it.
Li et al.\ (2011) \cite{li2011tortuosity} instead counted sharp bends along the main arteries of
consecutive patients with chest pain.
Tortuous patients were again less likely to have CAD, although the association was borderline, and
among patients with CAD the rate of adverse events at follow-up did not differ with tortuosity.
Most recently, Tello Ayala et al.\ (2026) \cite{telloayala2026tortuosity} segmented the right
coronary artery automatically in a large archive of angiograms from referred patients and computed
the turning angle at each point along its centerline.
Greater average tortuosity was associated with CAD, and a transformer that read the full profile of
turning angles identified CAD better than the average alone.
Depending on how tortuosity is measured, the same artery could therefore be read as a sign of less
disease or of more.

The geometric studies reviewed here examined only the right coronary artery or the left
tree alone \cite{telloayala2026tortuosity,garcha2025sensitivity,zhang2024curvature}.
To determine the full potential of the coronary tree as a risk marker, this must be extended.
Collectively, these findings suggest that an inconsistent measure
\cite{cui2017bifurcation,kashyap2022tortuosity} and a partial view of the tree may each hide where
geometry relates to disease.

Using geometry in routine risk assessment relies on automated segmentation of the coronary tree from
CCTA.
Bransby et al.\ (2026) \cite{bransby2026imagecasx} annotated a large public CCTA dataset anew and
benchmarked several deep learning segmentation methods on it.
They compared each method with the agreement between expert annotators of the same scans.
The best method came close to this agreement but ultimately concluded that automated
segmentation does not yet match human annotators.
